\scanpage{155}\mspage{10\makebox[0pt][r]{\tikz[baseline=0pt,overlay]{\draw (0.16em,0.70em) .. controls (0.16em,-0.37em) and (-1.16em,-0.39em) .. (-1.16em,0.44em);}}}
\unsure{généralisées|généralisés} \deleted{et plu\unsure{s}}\inserted{et le corps de classe ``\unsure{géométrique|géomrique|géomérique}''} (voir le livre de Serre), et plus tard sa théorie ``géométrique'' du corps de classe \uline{local} à coups de groupes \unsure{proalgébriques|pro.algébriques} (voir son article sur le sujet), m'avait fait réfléchir à des formulations cohomologiques de ces résultats\inserted{ et d'autres}, qui devaient être notamment de nature \inserted{``géom'' telle} que les résultats ``arithmétiques'' sur un corps de base (ou résiduel) $k$ \inserted{général (fini, p.ex)} se ramènent immédiatement par \unsure{descente|descent} au cas ``géométrique'' du corps de base $\bar{k}$. J'ai eu un échange de lettres avec Serre -- je ne sais si j'en retrouverai des copies, mais je me rappelle qu'il s'agissait de projets de théories assez ambitieux, \inserted{à coups de résidus généralisés, Jacobiennes gén. locales etc\unsure{,|;}} dans au moins trois directions différentes. Mais je ne suis jamais parvenu, malgré de nombreux essais, à accrocher une bonne volonté à un de ces programmes. Voici quelques mots sur eux :

C) Situation $X$ de type fini sur \uline{corps} $k$,

\noindent\begin{tabular}{@{}|p{0.94\linewidth}@{}}
Construction d'un complexe de jacobiennes généralisées $J_{*X/k}$ (de long. égale à $\dim X$)
\end{tabular}
% archive sheet number: 155

